% Generated draft: requires mathematical and source review.
\documentclass[11pt]{article}
\usepackage[margin=25mm]{geometry}
\usepackage{fontspec}
\setmainfont{Latin Modern Roman}
\usepackage{amsmath,amssymb,mathtools,unicode-math}
\setmathfont{Latin Modern Math}
\usepackage{enumitem,needspace,xurl,hyperref}
\setlength{\parindent}{0pt}
\setlength{\parskip}{5pt}
\setlength{\emergencystretch}{3em}
\newcommand{\N}{\mathbb{N}}
\newcommand{\Z}{\mathbb{Z}}
\newcommand{\Q}{\mathbb{Q}}
\newcommand{\R}{\mathbb{R}}
\newcommand{\C}{\mathbb{C}}
\newcommand{\F}{\mathbb{F}}
\newcommand{\A}{\mathbb{A}}
\newcommand{\PP}{\mathbb{P}}
\newcommand{\E}{\mathbb{E}}
\newcommand{\eps}{\varepsilon}
\newcommand{\dd}{\,\mathrm{d}}
\begin{document}
\section*{UnsolvedMath ID: 10400022}
\textbf{Source problem code:} catalogue entry
\subsection{Definitions and mathematical statement}
A torus knot \(K=T(p,q)\subset S^3\) is a simple closed curve of primitive slope \((p,q)\) on a standard embedded torus, with \(p,q\) coprime. The source asks to calculate a suitably defined Chern--Simons quantity \(\operatorname{CS}(S^3\setminus K)\) and a suitably branched logarithmic colored-Jones limit, extracted as
\[
 \lim_{N\to\infty}\frac{\operatorname{Log} J_N(K)}{N}.
\]
Giving the appropriate Chern--Simons definition and choosing the logarithm branch are expressly parts of the task.

For a knot complement, a Chern--Simons quantity generally requires a connection/representation, gauge group, normalization and boundary/peripheral data. The open complement and the compact exterior \(S^3\setminus\operatorname{int}\nu K\) require compatible boundary treatment. No canonical hyperbolic structure is assigned to a torus-knot complement.

The usual colored Jones polynomial \(J_N(K;q)\) needs a color convention and unknot normalization. The excerpt omits the evaluation defining the numerical \(J_N(K)\). A root such as \(q=e^{2\pi i/N}\) is a familiar possible context, but no specialization is inserted without the preceding definitions. Nonzero values and a coherent logarithm choice are required: changing branches by \(2\pi i k_N\) can change the normalized limit. Knot orientation/mirror, torus parameters and framing conventions must also be fixed.

The registry labels this solved, while the summary says ``Classification: PARTIAL-PROGRESS'', reports some Chern--Simons computations and an understood Jones limit, and gives no formulas or exact theorem references. The discrepancy remains unresolved rather than being treated as verification of a complete all-torus-knot answer.
\subsection{Short English statement}
For torus knots, define and compute the relevant Chern--Simons quantity and normalized logarithmic colored-Jones limit with an appropriate branch. The numerical specialization, boundary data and normalizations are missing, and the supplied solved/partial status claims conflict.
\subsection{Sources}
\begin{itemize}
\item Local registry: \url{data/unsolvedmath/problems.json}, permanent UnsolvedMath ID \(10400022\), source code catalogue entry; supplied statement, background, and research summary, updated 2026-08-06.
\item Ohtsuki, \emph{Problems on invariants of knots and 3-manifolds} (2002), Problem 1.22, attributed to H. Murakami, PDF page 25: \url{https://msp.org/gtm/2002/04/gtm-2002-04-024s.pdf}. Supplied PDF-font-structure extraction accessed 2026-07-29; provenance distinguishes 231 numbered items from surrounding references.
\item No separate Chern--Simons formula, Jones-asymptotic theorem or resolving publication is supplied.
\end{itemize}
\end{document}
